% Targeted replacements for the manuscript supplied on 2026-10-03.
% This is an insertion file, not a replacement for the entire source document.
% Preserve the original manuscript and its commented historical sections.
% Native-helper and agent progression values verified against the CECSL audit
% received 2026-10-03, which checked all 200 saved cases and frozen evidence.
% Artifact: equi-agent/outputs/audits/manuscript_received_20261003/.
% ZIP SHA256: 786045e67ce9d5a22b03f0ece073f1c81edebf16438602b64866d785e3cb94dc
% Standalone progression LLM rows remain from the supplied CECSL report;
% this archive did not re-audit their prediction files.

% ABSTRACT: replace the old GDP detection/progression superiority sentence.
% Do not substitute progression numbers into the separate detection result.
In an exploratory Harvard GDP progression analysis using baseline RNFLT maps
and visual-field measurements, RetinAgent increased positive-class F1 from
0.591 to 0.644 on the pointwise TD endpoint without a significance cutoff;
the gain did not extend consistently to the other progression endpoints.

% INTRODUCTION: replace the sentence containing 0.789/0.695 and 0.684/0.600.
On the Harvard GDP pointwise TD progression endpoint without a significance
cutoff, RetinAgent increased positive-class F1 from 0.591 for the native helper
to 0.644, and macro-F1 from 0.722 to 0.748. These results are exploratory:
the revised prompt was evaluated on the same cohort used to inspect earlier
agent errors, and performance gains did not generalize across all six endpoints.

% RESULTS 2.4: replace the progression paragraph and progression table only.
% The separate GDP detection paragraph still needs its canonical result audit.
Progression forecasting was evaluated from baseline RNFLT maps and 52-point
visual-field measurements, not from follow-up observations supplied at inference.
All methods were evaluated on the same 200 cases for each of six released GDP
endpoints. For \texttt{td\_pointwise\_no\_p\_cut}, RetinAgent achieved a
positive-class F1 of 0.6441 and a macro-F1 of 0.7476, compared with 0.5905 and
0.7224 for the native helper (Table~\ref{tab:gdp_progression_prediction}).
It detected 38 of 60 progressing cases, compared with 31 for the helper, while
false positives increased from 14 to 20 among 140 non-progressing cases.
Worst-group positive-class F1 increased from 0.5818 to 0.6102 among eligible
sex and age subgroups. No race comparison met the minimum support rule.
RetinAgent had
higher macro-F1 than the standalone language models in this comparison, but
GPT-5.1 and Claude Haiku 4.5 had slightly higher positive-class F1 and higher
balanced accuracy. Results varied substantially by endpoint: positive-class
F1 improved on the two pointwise TD endpoints, while no true positives were
identified by RetinAgent for MD, VFI or either rapid-MD endpoint
(Table~\ref{tab:gdp_progression_endpoint_summary}). Prompt revision followed
inspection of earlier errors on this cohort; independent evaluation is required
to establish generalization of the revised workflow.

\begin{table}[t]
\centering
\caption{Exploratory GDP progression forecasting for
\texttt{td\_pointwise\_no\_p\_cut}. All rows use 200 cases, including 60
progressing cases. Positive-class and macro-F1 are distinguished explicitly.
Worst-group F1 is positive-class F1 over race, sex and age groups with at least
20 positive and 20 negative cases and at least two eligible groups per
attribute. Only sex and age contributed eligible groups in this cohort;
race did not. RetinAgent's forced labels include cases flagged for review.
Reliability estimates use development-only out-of-fold predictions, but
prompt revision followed inspection of earlier test-cohort errors.}
\label{tab:gdp_progression_prediction}
\footnotesize
\begin{tabular}{@{}lcccccc@{}}
\toprule
Method & Positive F1 & Macro-F1 & Worst-group F1 & Sensitivity & Specificity & Balanced acc. \\
\midrule
GDP-native helper & 0.5905 & 0.7224 & 0.5818 & 0.5167 & 0.9000 & 0.7083 \\
GPT-5.1 & 0.6490 & 0.7181 & 0.5479 & 0.8167 & 0.7000 & 0.7583 \\
GPT-5.4 & 0.6420 & 0.6991 & 0.5679 & \textbf{0.8667} & 0.6429 & 0.7548 \\
GPT-5.6-luna & 0.2571 & 0.5498 & 0.1212 & 0.1500 & \textbf{0.9929} & 0.5714 \\
Claude Haiku 4.5 & \textbf{0.6494} & 0.7149 & 0.5479 & 0.8333 & 0.6857 & \textbf{0.7595} \\
RetinAgent & 0.6441 & \textbf{0.7476} & \textbf{0.6102} & 0.6333 & 0.8571 & 0.7452 \\
\bottomrule
\end{tabular}
\end{table}

\begin{table}[t]
\centering
\caption{Endpoint-specific performance of the GDP-native helper and revised
staged RetinAgent on the same 200 cases. This summary includes all six endpoints;
it is not an independently held-out evaluation of prompt revision.}
\label{tab:gdp_progression_endpoint_summary}
\footnotesize
\begin{tabular}{@{}lrrrrr@{}}
\toprule
Endpoint & Positives & \multicolumn{2}{c}{Positive F1} & \multicolumn{2}{c}{Balanced acc.} \\
 & & Helper & Agent & Helper & Agent \\
\midrule
MD & 18 & 0.0000 & 0.0000 & 0.4835 & 0.4368 \\
VFI & 19 & 0.0000 & 0.0000 & 0.4890 & 0.4448 \\
Pointwise TD & 18 & 0.1481 & 0.2174 & 0.5363 & 0.5757 \\
Rapid MD & 4 & 0.0000 & 0.0000 & 0.4923 & 0.4541 \\
Rapid MD, no p-cut & 6 & 0.0000 & 0.0000 & 0.4948 & 0.4510 \\
Pointwise TD, no p-cut & 60 & 0.5905 & 0.6441 & 0.7083 & 0.7452 \\
\bottomrule
\end{tabular}
\end{table}

% DISCUSSION paragraph 1: replace the broad GDP progression superiority claim.
In the revised GDP analysis, progression gains were endpoint-specific rather
than general: RetinAgent improved positive-class F1 over the native helper for
pointwise TD progression, with and without a significance cutoff, but did not
detect positive cases for the four remaining endpoints. On the no-cutoff TD
endpoint, improved sensitivity came with additional false positives. Because
prompt development included inspection of previous errors on the same cohort,
these results should be regarded as exploratory rather than confirmatory.

% METHODS, datasets: replace the FairVision split-size placeholder.
For each FairVision task, the full training, validation and test partitions
contained 6,000, 1,000 and 3,000 samples, respectively. The reduced agent
evaluation cohorts were drawn from the test partitions and were not used to
fit the foundation-model probes. The recovered glaucoma cohort contained
124 positive and 126 negative cases; AMD and DR locked manifests each
contained 125 positive and 125 negative cases.
% The sampling seed and canonical AMD agent coverage remain to be verified.

% METHODS, GDP: replace the claim that longitudinal visits were model inputs.
For progression forecasting, the released GDP data provide baseline RNFLT
maps, baseline visual-field values and six supplied progression labels for
500 patients. We used 300 development cases and 200 evaluation cases.
Five-fold development cross-validation generated out-of-fold reliability
estimates, with 240 fitting cases and 60 held-out development cases per fold.
Each final native helper was fitted on all 300 development cases for the fixed
training schedule, with a decision threshold of 0.5. No test cases contributed
to these reliability estimates. The helper was a supervised adaptation using
the released GDP model and dataset components, not a reproduction of the
paper's complete semi-supervised pseudo-supervisor experiment.

The primary endpoint was the released \texttt{progression[5]} label, named
\texttt{td\_pointwise\_no\_p\_cut}. The original paper describes pointwise TD
progression using deterioration at at least three visual-field locations;
the released index names this variant as omitting a p-value cutoff
\cite{luo2023harvard}. We used the supplied binary labels without recomputing
slopes or significance tests. Follow-up visits were not provided to the models
at inference. The other released indices correspond to MD, VFI, pointwise TD,
rapid MD, and rapid MD without a p-value cutoff.
% Source: https://arxiv.org/html/2308.13411v1, section 3.2;
% ../Harvard-GDP/README.md, released progression index mapping.
% Exact temporal horizon and label-generation details not supplied in the
% release must not be invented from the endpoint names.

% METHODS, reliability: replaces the fixed-before-test assertion and the
% one-subgroup shrinkage equation. This describes the historical FairVision
% score table. GDP uses the separate development-only construction below.
The reliability score was a manually specified composite, informed by
validation-prior ranking comparisons rather than fitted as a calibrated
probability of correctness. The largest coefficient was assigned to the
false-negative rate to emphasize missed disease; the false-positive term
penalized excessive referrals, while calibration error, discrimination and
F1 supplied complementary performance summaries. ECE used ten equal-width
probability bins on $[0,1]$, weighted by bin frequency. We do not claim that
the coefficients were selected before any test-result inspection. The
coefficient-sensitivity analysis therefore evaluates a heuristic score, not
a prespecified or optimized clinical utility function.

For a patient's age, race and sex strata, marginal risks were combined as
\begin{equation}
R_{\mathrm{local}} =
\frac{n_{\mathrm{age}}R_{\mathrm{age}}+
      n_{\mathrm{race}}R_{\mathrm{race}}+
      n_{\mathrm{sex}}R_{\mathrm{sex}}}
     {n_{\mathrm{age}}+n_{\mathrm{race}}+n_{\mathrm{sex}}},
\qquad
\widetilde R = \frac{n_{\cap}}{n_{\cap}+50}R_{\mathrm{local}}
 +\frac{50}{n_{\cap}+50}R_{\mathrm{global}}.
\end{equation}
The supplied trust score was $1-\widetilde R$. This is neither the worst
marginal risk nor a directly estimated intersectional performance metric.
Intersectional support $n_{\cap}$ controls shrinkage only. Undefined marginal
metrics fall back to their global counterparts. In the historical FairVision
table, performance priors came from validation predictions, whereas support
counts came from test-cohort demographic covariates. These two data sources
must not be described as a single validation-only estimate.

For the staged GDP experiment, both reliability metrics and support counts
came exclusively from the 300 development out-of-fold predictions. Marginal
groups required at least 20 positives and 20 negatives to contribute to the
local score; otherwise the global estimate was used. Reliability age groups
were younger than 40, 40--59 and 60 years or older. Evaluation age groups were
retained from the locked manifests and were not substituted by reliability
age groups.
% FairVision's lookup code maps age at 40/60 but the stored prior builder
% labels groups at 50/70. Resolve historical run provenance before inserting
% any uniform FairVision reliability age-bin claim.

% METHODS, staged GDP generation and errors: does not describe FairVision.
The staged GDP agent used the GPT-5.1 deployment with requested temperature
0.2 and a maximum completion budget of 8,000 tokens per stage. Five stages
were run per case: a Bio-Profiler, structural and functional specialists, an
evidence-ablation audit, and a final orchestrator. All six endpoint labels
were produced in the final response. The final and counterfactual stages
used explicit JSON schemas. Malformed or incomplete responses were retried
up to three total attempts with five seconds between attempts. Persistent
failures remained incomplete rather than being imputed as negative; the
main-table collector required 200 complete cases. Refusals and non-retryable
API errors stopped the call. Valid predictions flagged for clinician review
were retained in forced-classification metrics. Review flags were agent
judgements, not decisions made by an additional numeric escalation threshold.

For staged GDP, the five audit scenarios were full evidence, removal of
helper predictions, removal of structural interpretation, removal of
functional interpretation, and removal of demographic reliability context.
These scenarios were generated within one counterfactual-stage call. They
were not five randomly reassigned demographic profiles, independent votes,
or five complete executions of the pipeline.

% SUPPLEMENTARY METHODS AND RESULTS: offline sensitivity, not live ablation.
\subsection*{Offline coefficient sensitivity}
We reconstructed 486 saved FairVision foundation-model risk scores, covering
nine model--modality configurations, three tasks and 18 demographic
combinations per task. The reconstructed canonical scores matched the stored
scores to numerical precision. The analysis compared 21 settings: the
canonical weights, equal weights, the historical FNR/FPR-only score, removal
of each component, separate 20\% increases and decreases of each coefficient,
shrinkage constants of 25 and 100, and global-only scoring. Perturbed
coefficient vectors were renormalized to sum to one. We reported changes in
the minimum-risk model and absolute trust-score shifts; ties were retained.
Test labels were not used to choose or evaluate coefficient settings.

Across the ten single-coefficient 20\% perturbations, the minimum-risk model
was unchanged in all 18 glaucoma combinations. AMD changed in up to two
combinations, whereas DR changed in up to 12. Equal weighting changed the
minimum-risk model in zero glaucoma, three AMD and 15 DR combinations. These
findings show task-dependent ranking sensitivity. They do not establish
that live agent predictions or diagnostic performance are insensitive to the
weights; that requires a separate matched agent ablation.

% OPTIONAL SUPPLEMENTARY ANALYSIS: this may replace, but cannot be used to
% substantiate, the manuscript's live-agent counterfactual/escalation claim.
\subsection*{Separate deterministic metadata analysis}
A separate deterministic reliability-weighted fusion experiment used the
full 3,000-case test split for each FairVision task. It generated 54,000
single-attribute substitutions across 9,000 case--task records by replacing
race, ethnicity, sex or age-group values with alternative observed categories.
Diagnostic labels changed in 0.374\% of substitutions and escalation flags
in 5.811\%. This analysis used a different reliability weighting scheme
and a deterministic escalation policy; it was not a live RetinAgent
counterfactual experiment and cannot establish demographic invariance of
the live pipeline.
